\section{A uniform global bound for the finite evolutions}
\label{sec:global}

The estimates before $t_1$ give both families of finite evolutions a
common starting point for continuation.  We now prove that they exist
for all time, with a uniform bound on the part of the density that
contains the bath, the hot component, and the retired jets.  The
characteristic comparisons retain the color decomposition used in
\cite[Sections~5--6]{OpenAIPeriodic}.  The additional control comes from
the bath: the resulting small entropy distance of $F$ from $\mu$ supplies a
positive collision frequency and makes the nonlinear gain uniformly
bounded after averaging along free flights.
Entropy smallness and characteristic estimates also underlie the
bounded-amplitude theory of
\cite[Section~3]{DuanHuangWangYang2017}. Here the cold components remain
separate until retirement, while Lemma~\ref{lem:gain-small-mass}
controls the gain using a majorant singular on the velocity axis.

Recall from \eqref{con:late-colors} that, for $t\ge t_1$,
\begin{equation}\label{global:decomposition}
 F=c+w,\qquad c=\sum_{i\text{ active}}c_i,
 \qquad w(t_1)=u(t_1)+h(t_1).
\end{equation}
Between retirement times the equation for $w$ is
\begin{equation}\label{global:w-equation}
 (\D+\nu_F^b)w=P^b(w,w)+2P^b(c,w)+\mathcal S_c,
 \qquad
 \mathcal S_c=\sum_{i\ne i',\ i,i'\text{ active}}P^b(c_i,c_{i'}).
\end{equation}
The sum is ordered.  At $L_s=s^{-1/2}$ we add the jets of scale $s$ to
$w$ and remove them from $c$, so that $F$ has no jump.  All bounds below
hold for the characteristic representatives, at each time section and
almost everywhere in $(x,v)$, as in the finite evolution construction.

The bath bound in \cref{early:bath} and
\cref{early:continuation}, with $\zeta<1$, give a constant $B_0$ such that
\begin{equation}\label{global:initial-bound}
 \|u(t)+h(t)\|_H\le B_0\quad(0\le t\le t_1).
\end{equation}
Here $B_0$ is independent of the hole, the bath cutoff, the smallest
scale, and the approximation.  Fix provisionally $B>B_0+2$ and consider
an interval of existence beginning at $t_1$ on which
\begin{equation}\label{global:bootstrap}
 0\le w(t,x,v)\le B H(v).
\end{equation}
We will improve this bound by a constant independent of $B$, after
making the entropy error and the remaining small parameters sufficiently
small depending on $B$.

\subsection{Averaging and a positive collision frequency}

The spatial entropy estimate is an $L^1_x$ bound, whereas our
continuation argument needs estimates at almost every spatial point.
Velocity integration along flights makes this passage possible.  We
first record the elementary averaging estimate in a form that also
applies over arbitrarily long time intervals.  Write
$M_H=\int H(v)\dd v<\infty$.

\begin{lemma}[Averaging along free flights]\label{global:flight-average}
Let $A:[0,\infty)\times(\R^3/\Lambda)\to[0,\infty)$ be measurable,
write $A_s=A(s,\cdot)$, and suppose
\[
 \sup_s\|A_s\|_\infty\le A_*,\qquad
 \sup_s\|A_s\|_1\le a_*.
\]
For a fixed bounded velocity set $V$ and $\rho>0$,
\begin{equation}\label{global:ball-average}
 \int_V A_s(x-\tau v)\dd v
 \le C_{V,\rho}a_*
 \qquad(\tau\ge\rho),
\end{equation}
uniformly in $x,s,\tau$.  If $A_*$ is fixed, the following quantities
tend to zero as $a_*\to0$:
\begin{align}
 &\sup_{t\ge t_1,x}
 \int_{t_1}^t e^{-\kappa(t-s)}
       \int H(v)A_s(x-(t-s)v)\dd v\dd s,
       \label{global:damped-average}\\
 &\sup_{t\ge t_1,x}
 \int_0^{t_1}\int H(v)A_s(x-(t-s)v)\dd v\dd s,
       \label{global:early-average}
\end{align}
where $\kappa>0$ and $t_1>0$ are fixed.  The supremum in $x$ is essential.
\end{lemma}

\begin{proof}
Contain $V$ in a ball of radius $R$.  The change of variables
$z=x-\tau v$ gives a factor $\tau^{-3}$.  The image ball meets at most
$C(1+\tau R)^3$ translates of a fixed bounded fundamental cell of
$\Lambda$.  Periodicity therefore gives
\[
 \int_V A_s(x-\tau v)\dd v
 \le C\tau^{-3}(1+\tau R)^3\|A_s\|_1,
\]
which proves \eqref{global:ball-average}.  In particular, its constant
does not grow when $\tau$ tends to infinity.

For $L\ge1$ let
\[
 H_L(v)=\min\{H(v),L\}\ind_{\{|v|\le L\}},\qquad
 \delta_L=\int(H-H_L)\dd v.
\]
Then $\delta_L\to0$, and \eqref{global:ball-average} implies, for
$\tau\ge\rho$,
\begin{equation}\label{global:weighted-average}
 \int H(v)A_s(x-\tau v)\dd v
 \le A_*\delta_L+C_{L,\rho}a_*.
\end{equation}
For $0\le\tau<\rho$, the same integral is at most $A_*M_H$.
Consequently the expression in \eqref{global:damped-average} is at most
\[
 \rho A_*M_H+\kappa^{-1}
       (A_*\delta_L+C_{L,\rho}a_*),
\]
and that in \eqref{global:early-average} is at most
\[
 \rho A_*M_H+t_1(A_*\delta_L+C_{L,\rho}a_*).
\]
Choose first $\rho$ small, then $L$ large, and finally $a_*$ small.
Both upper bounds tend to zero in this order.  The same proof works
when the time integrals are restricted to an interval of existence.
\end{proof}

Recall the entropy quantities from \eqref{con:entropy-distance}:
\begin{equation}\label{global:entropy-distance}
 \begin{aligned}
 D_F(t,x)&=\int\langle v\rangle|F(t,x,v)-\mu(v)|\dd v,
 &\|D_F(t)\|_1&\le\varepsilon_E,\\
 D_-(t,x)&=\int\langle v\rangle(\mu(v)-F(t,x,v))_+\dd v.
 \end{aligned}
\end{equation}
The second quantity is bounded pointwise by
$\int\langle v\rangle\mu(v)\dd v$, independently of the bootstrap.
The first has a pointwise bound
\begin{equation}\label{global:moment-bound}
 D_F(t,x)\le C(B),\qquad t\ge t_1/2,
\end{equation}
as long as \eqref{global:bootstrap} holds after $t_1$.  Indeed, before
$t_1$ use \eqref{global:initial-bound}; after $t_1$ use $w\le BH$.
At every point there is at most one active main jet, whose mass and
first moment are bounded by $C/(t+s)\le2C/t_1$ by
\cref{cold:core-amplitude}.  The other jets are controlled by
\cref{cold:nonmain}.  These bounds prove \eqref{global:moment-bound}
uniformly in the number of retained scales.

\begin{lemma}[A collision frequency independent of the bootstrap size]
\label{global:frequency}
There is a constant $\kappa>0$, depending only on the fixed Maxwellian
$\mu$ and $t_1$, with the following property.  For each fixed $B>B_0+2$,
if $\varepsilon_E$ is sufficiently small depending on $B$, then every
finite evolution satisfying \eqref{global:bootstrap} obeys
\begin{equation}\label{global:frequency-bound}
 \nu_F^b(t,x,v)\ge\kappa\qquad(t\ge t_1).
\end{equation}
The constant is uniform for $b\ge1$ and for all approximation indices.
\end{lemma}

\begin{proof}
Fix two closed velocity balls $V_1,V_2$ of positive radius with
$\dist(V_1,V_2)>3$, and put $V=V_1\cup V_2$.  All constants until the
choice of the entropy tolerance may depend on these fixed balls.
Since $b\ge1$ and $\mu$ is a Maxwellian,
\[
 P^b(\mu,\mu)(v)=\mu(v)\nu_\mu^b(v)
 \ge \mu(v)\nu_\mu^1(v)\ge m_0>0
 \qquad(v\in V).
\]
At a fixed $(t,x)$ write $d=(\mu-F)_+$.  Positivity of the gain and
$F\ge\mu-d$ imply
\[
 P^b(F,F)\ge P^b(\mu-d,\mu-d)
 \ge P^b(\mu,\mu)-2P(\mu,d).
\]
By the plane formula \eqref{ker:carleman}, the Gaussian integral over
each plane is bounded, so
\[
 P(\mu,d)(v)\le C\int\frac{d(p)}{|p-v|}\dd p.
\]
For $|p-v|<r$, use $d\le\mu\le1$ to bound this integral by $Cr^2$;
for $|p-v|\ge r$, it is at most $Cr^{-1}D_F(t,x)$.  Choose $r>0$
once so that the first contribution to $2P(\mu,d)$ is at most $m_0/2$.
With $m_*=m_0/2$ and a fixed constant $C$, we obtain
\begin{equation}\label{global:gain-lower}
 P^b(F,F)(t,x,v)\ge m_*-CD_F(t,x),\qquad
 \nu_F^b(t,x,v)\le C+CD_F(t,x),\qquad v\in V.
\end{equation}
The second inequality follows directly from
$\min\{b,|v-p|\}\le |v-p|\le C\langle p\rangle$ on $V$.

Let $\Delta=t_1/2$, fix $t\ge t_1$, and write
$X_\tau=x-(t-\tau)v$.  The Duhamel formula for the total density,
whose value is continuous through retirements, gives
\[
 F(t,x,v)\ge
 \int_{t-\Delta}^{t-\Delta/2}
 e^{-\int_s^t\nu_F^b(\tau,X_\tau,v)\dd\tau}
 P^b(F,F)(s,X_s,v)\dd s.
\]
Using \eqref{global:gain-lower}, $e^{-a}\ge1-a$ for $a\ge0$, and the
upper bound $1$ for a survival factor, the integrand is bounded below by
\begin{equation}\label{global:survival-lower}
 m_*e^{-C\Delta}
 \left(1-C\int_s^t D_F(\tau,X_\tau)\dd\tau\right)
 -CD_F(s,X_s).
\end{equation}
This inequality remains valid when its right side is negative.

Integrate \eqref{global:survival-lower} over $s$ and over either ball
$V_k$.  The positive constant contributes
\[
 a_k:=|V_k|(\Delta/2)m_*e^{-C\Delta}>0.
\]
Choose $0<\rho<\Delta/2$.  For $t-\tau\ge\rho$,
\eqref{global:ball-average} and \eqref{global:entropy-distance} give
\begin{equation}\label{global:entropy-flight}
 \int_{V_k}D_F(\tau,x-(t-\tau)v)\dd v
 \le C_{V_k,\rho}\varepsilon_E.
\end{equation}
The term $D_F(s,X_s)$ always has $t-s\ge\Delta/2$.  Its integrated
cost is therefore $C_\rho\varepsilon_E$.  For the double time integral
in \eqref{global:survival-lower}, the portion with $t-\tau\ge\rho$ has
the same type of bound, whereas the remaining portion costs at most
$C(B)\rho$ by \eqref{global:moment-bound}.  Thus
\[
 \int_{V_k}F(t,x,v)\dd v
 \ge a_k-C_\rho\varepsilon_E-C(B)\rho.
\]
After $B$ has been fixed, choose $\rho$ to make the last term at most
$\min(a_1,a_2)/4$, and then choose $\varepsilon_E$ to make the middle
term at most the same quantity.  It follows that
\[
 \int_{V_k}F(t,x,v)\dd v\ge a_k/2\qquad(k=1,2).
\]
For any velocity $v$, one of the balls has distance at least $1$ from
$v$: otherwise their mutual distance would be less than $2$.
On that ball $\min\{b,|v-p|\}\ge1$.  Consequently
\[
 \nu_F^b(t,x,v)
 =2\pi\int\min\{b,|v-p|\}F(t,x,p)\dd p
 \ge\pi\min(a_1,a_2)=:\kappa.
\]
Only the entropy tolerance, and not this value of $\kappa$, depended
on $B$.
\end{proof}

The positive frequency gives exponential decay to every active jet.
Specifically, \eqref{cold:decay} yields
\begin{equation}\label{global:cold-decay}
 c_i(t,x,v)\le e^{-\kappa(t-t_1)}\mathcal U_i(t,x,v),
 \qquad t_1\le t\le L_{s_i}.
\end{equation}
At a point in the active tube union $\mathcal O(t)$ there is at most
one main color.  Its sheet amplitude is at most $C/t_1$; outside the
union all colors are nonmain.  Applying \cref{lem:sheet-bound} and
\cref{cold:nonmain}, and enlarging its small coefficient if necessary,
gives constants $C_3,C_4$ independent of $B$ such that
\begin{align}
 P(c,H)&\le C_3e^{-\kappa(t-t_1)}H
            \bigl(\ind_{\mathcal O(t)}+\theta_0\bigr),
              \label{global:cold-gain}\\
 \mathcal S_c&\le C_4\theta_0H,
              \label{global:cold-source}
\end{align}
where $\theta_0=\theta_0(s_0)\to0$ as $s_0\to0$.
For the second estimate, if $c_*$ is the main color and $r$ the sum of
the remaining colors, the ordered cross sum is at most
$2P(c_*,r)+P(r,r)$.  Here $r\le\theta_0e^{-4|v|^2}$ and the sheet bound
controls $P(c_*,H)$; if there is no main color use $c\le
\theta_0e^{-4|v|^2}$ and \cref{lem:hot-bound}.

At a retirement time, \cref{cold:retirement-size} and
\eqref{global:cold-decay} bound the increment of $w$ by $J_sH$, where
\begin{equation}\label{global:retirement}
 J_s=C s^{-C}e^{-\kappa(L_s-t_1)},\qquad
 \sum_sJ_s\le J_0(s_0):=
 \sum_{j\ge1}C s_j^{-C}e^{-\kappa(s_j^{-1/2}-t_1)}\longrightarrow0.
\end{equation}
This follows for the one-sided trace just before retirement, using
only the bootstrap up to that time.  To check the last limit explicitly,
the exponential dominates every power of $s$: for sufficiently small
$s$, $Cs^{-C}e^{-\kappa(s^{-1/2}-t_1)}\le s^2$.  The remaining geometric
sum is at most $s_0^2/(e^{2\eta}-1)$.  In particular the sum of all
retirement increments can be made small uniformly in $N$.

\subsection{Small excess mass and a bounded nonlinear gain}

The remaining obstruction is the quadratic term $P^b(w,w)$: estimating
it directly from $w\le BH$ would produce a constant of order $B^2$.
Instead we compare $w$ with the fixed Maxwellian.  Entropy bounds the
spatial integral of its positive excess, and the preceding flight
estimate converts that bound into small excess mass at each position.
The small-mass gain estimate then controls the quadratic term by a
constant chosen before $B$.

\begin{lemma}[Uniform nonlinear gain]\label{global:gain}
Let $B>B_0+2$ be fixed and suppose \eqref{global:bootstrap} holds on an
interval beginning at $t_1$.  After taking
$\varepsilon_E$, $\zeta$, and $s_0$ sufficiently small depending on $B$,
one has on that interval
\begin{equation}\label{global:gain-bound}
 P^b(w,w)\le C_5H,
 \qquad C_5:=\|P(\mu,\mu)\|_H+1.
\end{equation}
In particular, $C_5$ is independent of $B$, the approximation, and the
cap horizon.  The estimate holds also at the appropriate one-sided
sections of retirement times.
\end{lemma}

\begin{proof}
Choose the entropy tolerance small enough for
\cref{global:frequency}, so that the damping constant $\kappa$ and
\eqref{global:cold-gain}--\eqref{global:retirement} are available.
Define
\begin{equation}\label{global:excess}
 j=(w-\mu)_+,\qquad m_j(t,x)=\int j(t,x,v)\dd v.
\end{equation}
Since $w\le F$ and $j\le BH$,
\begin{equation}\label{global:excess-mass}
 0\le m_j\le BM_H,
 \qquad\|m_j(t)\|_1\le\varepsilon_E.
\end{equation}
Write $C_\mu=\|\mu\|_H<\infty$ and $A_B=2C_\mu+B$.
The inequality $w\le\mu+j$, bilinearity, and positivity give
\begin{equation}\label{global:gain-difference}
 P^b(w,w)-P^b(\mu,\mu)
 \le2P(\mu,j)+P(j,j)\le A_BP(H,j).
\end{equation}
The frequency deficit has the bound
\begin{equation}\label{global:loss-difference}
 (\nu_\mu^b-\nu_F^b)\mu
 \le2\pi\mu(v)\int|v-p|(\mu(p)-F(p))_+\dd p
 \le CH(v)D_-.
\end{equation}
Here $\mu(v)\langle v\rangle\le CH(v)$.

Between retirements, subtracting the stationary Maxwellian equation
from \eqref{global:w-equation} gives
\[
 (\D+\nu_F^b)(w-\mu)
 =P^b(w,w)-P^b(\mu,\mu)
   +(\nu_\mu^b-\nu_F^b)\mu+2P^b(c,w)+\mathcal S_c.
\]
The chain rule for the positive part on almost every free flight
therefore gives an upper comparison for $j$ with damping $\kappa$.
Using \cref{lem:gain-small-mass} in \eqref{global:gain-difference},
and then \eqref{global:loss-difference},
\eqref{global:cold-gain}, and \eqref{global:cold-source}, shows that
for each $\epsilon>0$ its positive source is at most $H(v)R(s,x)$,
where
\begin{equation}\label{global:excess-source}
 R(s,x)=\epsilon+C_{B,\epsilon}m_j(s,x)+CD_-(s,x)
           +C(B)\bigl(\ind_{\mathcal O(s)}(x)+\theta_0\bigr).
\end{equation}
Constants have been enlarged to absorb $A_B$.  At a retirement time,
the increase of $j$ is at most the increase of $w$, because
$(a+d)_+-a_+\le d$ for $d\ge0$.  Thus its increments are bounded by
the same $J_sH$ as in \eqref{global:retirement}.

We also need to estimate the excess already present in the bath at
$t_1$.  Put $j_u=(u-\mu)_+$ and $m_{j_u}=\int j_u\dd v$ for
$0\le s\le t_1$.  The bath starts below $\mu$, so $j_u(0)=0$;
\cref{early:bath} gives the fixed bound $j_u\le B_0H$.
Since $u\le F$,
\begin{equation}\label{global:bath-excess-mass}
 \|m_{j_u}(s)\|_\infty\le B_0M_H,
 \qquad\|m_{j_u}(s)\|_1\le\varepsilon_E.
\end{equation}
The bath equation is $(\D+\nu_F^b)u=P^b(u,u)$.  Subtracting the
Maxwellian equation and applying the same two difference estimates as
above gives, now discarding the nonnegative damping, a source bounded
by $H(v)R_u(s,x)$, where
\begin{equation}\label{global:bath-excess-source}
 R_u(s,x)=\epsilon+C_{B_0,\epsilon}m_{j_u}(s,x)+CD_-(s,x).
\end{equation}
The cold density appears only through the frequency deficit, which is
controlled by $D_-$.  No positive source involving a cold gain occurs
in this bath equation.

At $t_1$, the hot bound gives $j(t_1)\le j_u(t_1)+\zeta H$.
Continue the same backward free flight through $t_1$, and write
$X_s=x-(t-s)v$.  The two comparisons, together with all intervening
retirements, imply
\begin{equation}\label{global:excess-path}
 \frac{j(t,x,v)}{H(v)}
 \le\zeta+J_0
   +\int_0^{t_1}R_u(s,X_s)\dd s
   +\int_{t_1}^t e^{-\kappa(t-s)}R(s,X_s)\dd s.
\end{equation}
We have dropped damping from the initial and retirement contributions;
doing so also allows the bath comparison to run along that one flight
all the way back to zero.

We claim that the velocity integral of the right side times $H$ can
be made smaller than any prescribed number $\tau>0$, uniformly in
$(t,x)$ on the bootstrap interval.  The constant terms contribute at
most
\begin{equation}\label{global:constant-errors}
 M_H\left(\zeta+J_0+\epsilon(t_1+\kappa^{-1})
                      +\kappa^{-1}C(B)\theta_0\right).
\end{equation}
For the remaining terms, all the spatial functions in question have
a fixed pointwise upper bound once $B$ is fixed:
$m_j\le BM_H$, $m_{j_u}\le B_0M_H$,
$D_-\le\int\langle v\rangle\mu\dd v$, and
$\ind_{\mathcal O(s)}\le1$.  Their spatial integrals are respectively
bounded by $\varepsilon_E$, $\varepsilon_E$, $\varepsilon_E$, and
$Cs_0^{3/2}$, the last by \eqref{cold:tube-volume}.
Apply \cref{global:flight-average} to each term, with its coefficient
in \eqref{global:excess-source} or \eqref{global:bath-excess-source}.
For clarity, the choices are made in the following order.  First choose
$\epsilon$ so that its contribution to \eqref{global:constant-errors}
is less than $\tau/4$; all the coefficients $C_{B,\epsilon}$ and
$C_{B_0,\epsilon}$ are then fixed.  For each of the finitely many
averaged terms, choose the time deletion length $\rho$ and the
truncation $H_L$ in the proof of that lemma so that the endpoint and
velocity-tail errors, including its coefficient, have total less than
$\tau/4$.  Next make $\varepsilon_E$ and $s_0$ sufficiently small for
the remaining averaged errors to total less than $\tau/4$.  Finally
take $\zeta$ small and, if necessary, reduce $s_0$ further to make the
other terms in \eqref{global:constant-errors} less than $\tau/4$.
All these choices are uniform in the endpoint time.  Thus
\begin{equation}\label{global:small-excess}
 m_j(t,x)\le\tau.
\end{equation}
The same calculation on a pre-retirement section includes only the
jumps that have already occurred, so it does not require a post-jump
bootstrap hypothesis.

It remains to select the particular smallness threshold that controls
the gain.  Apply \cref{lem:gain-small-mass} once more, now with
$\epsilon'=1/(2A_B)$.  Write its constant as $C_{B,\epsilon'}$ and set
\[
 \tau=\frac{1}{2A_B\max\{1,C_{B,\epsilon'}\}}.
\]
Choose the parameters in the preceding argument to give
\eqref{global:small-excess} for this $\tau$.  Then
\[
 A_BP(H,j)\le
 A_BH\bigl(\epsilon'+C_{B,\epsilon'}m_j\bigr)\le H.
\]
Equation \eqref{global:gain-difference} and $P^b(\mu,\mu)\le
P(\mu,\mu)$ prove \eqref{global:gain-bound}.  The latter has finite
$H$ norm because $P(\mu,\mu)=\mu\nu_\mu$ is a Gaussian times a function
of at most linear growth.
\end{proof}

\subsection{Parameter choice and continuation}

We can now choose the bootstrap size and then realize all the
smallness requirements imposed above.  This order is essential: both
the damping rate and the ultimate gain bound have already been fixed
independently of the provisional value $B$.

\begin{proposition}[Uniform global continuation]\label{prop:global-control}
The construction parameters can be fixed so that every dormant finite
evolution and every cap solution with $0<T\le T_0$ continues globally.
There are constants $B<\infty$ and $\kappa>0$, independent of $N$,
$b_N\ge1$, and $T$, such that
\begin{equation}\label{global:uniform-bound}
 \|w(t)\|_H\le B,\qquad \nu_F^{b_N}(t,x,v)\ge\kappa,
 \qquad t\ge t_1.
\end{equation}
Every active jet satisfies \eqref{global:cold-decay} through its
retirement, the total retirement increment is bounded by
\eqref{global:retirement}, and the hot component satisfies the small
bound of \cref{early:continuation} on $[0,t_1]$.
\end{proposition}

\begin{proof}
On a bootstrap interval let $Z_w(t)=\|w(t)\|_H$.  The Duhamel formula
for \eqref{global:w-equation}, with all retirement increments included,
and \eqref{global:frequency-bound}, \eqref{global:cold-gain},
\eqref{global:cold-source}, and \eqref{global:gain-bound}, give
\begin{align*}
 Z_w(t)\le {}&e^{-\kappa(t-t_1)}B_0
  +\frac{C_5+C_4\theta_0}{\kappa}+J_0\\
 &+2C_3(1+\theta_0)
     \int_{t_1}^t e^{-\kappa(t-s)}e^{-\kappa(s-t_1)}Z_w(s)\dd s.
\end{align*}
Take $s_0$ small enough that $\theta_0,J_0\le1$ and discard the first
exponential in the last integrand.  We obtain
\begin{equation}\label{global:scalar-bootstrap}
 Z_w(t)\le B_0+\frac{C_4+C_5}{\kappa}+1
       +4C_3\int_{t_1}^t e^{-\kappa(s-t_1)}Z_w(s)\dd s.
\end{equation}
The integral version of Gronwall's inequality therefore gives
\begin{equation}\label{global:strict-improvement}
 Z_w(t)\le
 \left(B_0+\frac{C_4+C_5}{\kappa}+1\right)
       \exp\left(\frac{4C_3}{\kappa}\right)=:B_*.
\end{equation}
The constants defining $B_*$ were fixed before $B$.

Here is the complete order of parameter choices.  Fix the kernel and
jet parameters, the narrow width factor, the Maxwellian $\mu$, and the
short time $t_1$ as in the preceding sections; also fix the early
growth exponent $p$ there.  The bath estimate fixes $B_0$.
The proof of \cref{global:frequency} fixes $\kappa$ from $\mu$, $t_1$,
and the two chosen balls, without imposing its entropy tolerance yet.
The cold estimates then fix $C_3,C_4$ uniformly for sufficiently small
$s_0$, while $C_5=\|P(\mu,\mu)\|_H+1$ is already fixed.  Choose
\[
 B>\max\{B_0+2,B_*+3\}.
\]
For this $B$, prescribe the entropy tolerance required by
\cref{global:frequency}, and the possibly smaller entropy tolerance,
hot tolerance $\zeta$, and upper bounds on $s_0$ required by
\cref{global:gain}.  Include $\theta_0,J_0\le1$ among these requirements.

To realize the entropy tolerance, take the hole radius $\delta$ small
and the bath velocity cutoff $R_\mu$ large.  Their contributions to
the initial relative entropy can then be made arbitrarily small.
With these parameters fixed, the jet and seed contributions become
arbitrarily small as $s_0\to0$, as proved in
\eqref{con:entropy-distance} and its derivation.  To realize the
prescribed $\zeta$, use \cref{early:continuation}: choose its auxiliary
time $t_b$ and horizon upper bound $T_0$ after $\delta,\zeta$, and then
take $s_0$ sufficiently small for its source estimate.  Decrease $s_0$
further to meet the entropy, separation, cold-gain, and retirement
requirements already listed.  Reducing $T_0$ when required by the
early estimates preserves all later estimates.  These are fixed
construction parameters; none depends on $N$, on the chosen cap solution,
or on its horizon $T\in(0,T_0]$.

For completeness, the strict estimate closes the bootstrap even at
retirement times.  At $t_1$ it starts by
\eqref{global:initial-bound}.  On an interval with no retirement,
the characteristic equations make the weighted norm continuous, and
\eqref{global:strict-improvement} bounds it by $B_*<B-3$ as long as
it is at most $B$.  Thus it cannot first reach $B$ there.  At a
retirement time $L_s$, apply the argument only up to the left-hand
section.  The cold comparison on that interval and
\eqref{global:retirement} give
\[
 Z_w(L_s+)\le Z_w(L_s-)+J_s\le B_*+1<B-2.
\]
This starts the bootstrap on the next interval without having assumed
its conclusion across the jump.  The total density and its entropy
bound are unchanged by this relabeling.

Finally, at each fixed approximation only finitely many scales occur.
The cold barriers bound the sum of the active colors in $H$ norm on
every finite interval, with constants allowed to depend on that
approximation.  The uniform bound on $w$ therefore bounds the total
$H$ norm on every such interval.  The cutoff local existence and
continuation argument from the construction excludes a finite maximal
existence time.  There are only finitely many retirements at that
approximation, after which $F=w$.  This proves global existence and
\eqref{global:uniform-bound}, with all the stated auxiliary estimates.
\end{proof}
